\section{Identifizierte Erfolgsfaktoren und Muster}
\label{sec:05-04_success-factors}

Auf Basis der triangulierten Auswertung der \Gls{usability}-Metriken, Skalenbewertungen und qualitativen Interviewdaten lassen sich mehrere übergreifende Erfolgsfaktoren und stabile Muster identifizieren. Diese treten unabhängig von einzelnen Aufgaben, Rollen oder Nutzungshäufigkeiten auf und beschreiben wiederkehrende Wirkmechanismen der neuen Wiki-Gestaltung.

Ein zentraler Erfolgsfaktor ist die klare und konsistente \acrlong{ia}. Dieses Muster wird sowohl durch hohe subjektive Orientierungsbewertungen als auch durch deutliche Effizienz- und Effektivitätsgewinne in den \Gls{usability}-Tests sowie reduzierte Klickzahlen gestützt. Qualitativ wurde die Struktur von der Mehrheit der Teilnehmenden als logisch nachvollziehbar beschrieben, insbesondere im Hinblick auf flache Navigationshierarchien, eindeutige Kategorienbezeichnungen und thematische Bündelungen. Die Übereinstimmung zwischen subjektiver Wahrnehmung und objektiven Leistungsdaten verdeutlicht, dass strukturelle Klarheit unmittelbar in messbare Nutzungsvorteile übersetzt wird.

Ein zweites wiederkehrendes Muster betrifft die Bedeutung von Einstiegs- und Übersichtsseiten als Orientierungsanker. Mehrere Teilnehmende beschrieben diese Seiten als zentralen Ausgangspunkt für Exploration und Wiedereinstieg, insbesondere bei seltener Nutzung oder komplexeren Informationsbedarfen. Dieses qualitative Muster korrespondiert mit besonders ausgeprägten Zeit- und Fehlerreduktionen bei Aufgaben mit hohem Such- und Orientierungsanteil. Einstiegsseiten fungieren damit nicht nur als subjektiv wahrgenommene Unterstützung, sondern als funktionales Element zur Reduktion kognitiver und navigationaler Aufwände.

Ein drittes Muster zeigt sich in der hohen Konsistenz zwischen subjektiver Bewertung und objektiver Nutzungseffizienz. Die Skalenbewertungen zur Orientierung, Rollenpassung und Nutzerzentrierung weisen durchgehend hohe Zustimmungswerte bei geringer Streuung auf und werden durch die Leistungskennzahlen der \Gls{usability}-Tests sowie die Klickdaten bestätigt. Abweichungen zwischen Wahrnehmung und messbarer Leistung treten lediglich punktuell auf, insbesondere im Bereich der Personen- und Ansprechpartnerlogik. Insgesamt deutet dieses Muster darauf hin, dass die Nutzenden ihre Gebrauchstauglichkeitserfahrungen weitgehend realistisch einschätzen.

Demgegenüber lässt sich als stabiles Negativmuster die eingeschränkte Verständlichkeit der Personen- und Ansprechpartnerlogik identifizieren. Unabhängig von Rolle oder Erfahrungsgrad äußerte eine Mehrheit der Teilnehmenden Kritik an der Platzierung und Struktur dieses Inhaltsbereichs. Dieses qualitative Muster spiegelt sich in den quantitativen Daten wider, da die entsprechende Aufgabe im Vergleich zu allen übrigen Aufgaben keine relevanten Leistungsgewinne und keine Reduktion der Klickanzahl zeigte. Der Befund verdeutlicht, dass lokale Inkonsistenzen innerhalb einer insgesamt verbesserten \acrlong{ia} gezielt wahrgenommen werden und unmittelbare Auswirkungen auf Nutzungseffizienz, Fehlerraten und Navigationsaufwand haben.

Ein weiteres übergreifendes Muster betrifft die Akzeptanz generischer statt stark personalisierter Strukturen. Obwohl vereinzelt der Wunsch nach stärkerer Rollendifferenzierung geäußert wurde, bewerteten nahezu alle Teilnehmenden die vorhandene Struktur als ausreichend passend für ihre jeweilige Rolle. Entscheidend ist weniger eine explizite Personalisierung als vielmehr die Verfügbarkeit vollständiger Inhalte, transparente Zuständigkeiten und eine gute Auffindbarkeit. Rollenadäquatheit entsteht damit implizit aus der Strukturqualität und nicht aus individuellen Anpassungsmechanismen.

Darüber hinaus zeigt sich ein konsistentes Muster in der Gewichtung funktionaler gegenüber rein visuellen Optimierungen. Verbesserungsvorschläge und Kritikpunkte beziehen sich überwiegend auf funktionale Aspekte wie Suchunterstützung, Auffindbarkeit von Ansprechpartnern und Navigationslogik, während visuelle Gestaltungselemente überwiegend positiv und nur selten als Optimierungsbedarf genannt wurden. Dieses Muster deutet darauf hin, dass aus Nutzendensicht funktionale Unterstützung einen höheren Einfluss auf wahrgenommene Gebrauchstauglichkeit und Akzeptanz besitzt als ästhetische Feinjustierungen.

Schließlich zeigt sich ein weiteres stabiles Muster in der impliziten Wahrnehmung von \Gls{uxstrategien}. Ein Großteil der Teilnehmenden schloss aus der Klarheit der Struktur, der reduzierten visuellen Gestaltung, konsistenter Benennung und der Orientierung an Nutzungsszenarien auf eine nutzerzentrierte Entwicklung. Der Einsatz formaler \Gls{ux}-Methoden wurde dabei weniger explizit erkannt, sondern primär über wahrnehmbare Gestaltungseigenschaften erschlossen. \Gls{ux} wird somit im Nutzungskontext nicht als Methode, sondern als Qualität des Nutzungserlebnisses wahrgenommen.

Zusammenfassend lassen sich als zentrale Erfolgsfaktoren insbesondere eine klare \acrlong{ia}, unterstützende Einstiegslogiken, konsistente Strukturierung sowie die hohe Übereinstimmung zwischen subjektiver Wahrnehmung und objektiver Nutzungseffizienz identifizieren. Gleichzeitig bestehen stabile Schwachstellen im Bereich der Personen- und Ansprechpartnerlogik sowie in funktionalen Unterstützungsmechanismen. Die identifizierten Muster verdeutlichen, welche Gestaltungsprinzipien unmittelbar zur Gebrauchstauglichkeit beitragen und welche strukturellen Aspekte weiterhin Optimierungspotenzial aufweisen. Diese Befunde bilden die Grundlage für die abschließende Betrachtung der Akzeptanzfaktoren im organisationalen Kontext im folgenden Kapitel.


